\scanpage{81}
En fait, les relations $-2<\alpha_0,\alpha_1<+2$ sont satisfaites automatiquement, si $\alpha_0,\alpha_1$ solutions d'équations $\tfrac12$ cyclotomiques, correspondant à $\nu_0,\nu_1\ne1,2$ i.e. $\nu_0,\nu_1>3$. Quitte à passer au système dual, on est donc amené à regarder les possibilités, où

\[
 3\leq\nu_0\leq\nu_1\leq5
\]

A priori, il y a \deleted{\unsure{$(3,4)(3,5$|$(3,3)(3,5)$}} $(3,3)(3,4)(3,5)$ (cas I II III) $(4,4),(4,5)$ et $(5,5)$. Or $(4,4)$ correspond à \deleted{\unsure{$(0,0$}} $\alpha_0=\alpha_1=0$, $\alpha_0+\alpha_1=0$, impossible, et $(4,5)$ à $\alpha_0=0,\alpha_1=\tfrac12(-1\pm\sqrt5)$. Mais \uline{OPS} $\alpha_1=\tfrac12(-1\underline{\underline{+}}\sqrt5)$ qui est $>0$, \deleted{on a gagné} impossible ! Enfin considérons le cas $5,5$ OPS $\alpha_0=\tfrac12(-1\underline{\underline{+}}\sqrt{\unsure{5|6}})>0$, donc il faut $\alpha_1<0$ donc $\alpha_{1}=\tfrac12(-1-\sqrt{\deleted{\unsure{5}}5})$ ($\alpha_0+\alpha_1=-1$ OK).

On a gagné !
% archive sheet number: 80
